\chapter{二分类学习的可学习性}
\label{chap:binary-learnability}

第~\ref{chap:statistical-learning-framework}章中，阈值0.4和0.6对应的假设都分对了四封训练邮件，
却对分数为0.5的新邮件给出不同答案。增加训练样本有机会区分它们：如果某个假设经常出错，
新样本就更容易落在它出错的地方。然而，算法会在观察样本后从整个假设空间中作出选择。
某个假设即使经常把新邮件分错，也可能碰巧分对全部训练邮件；只看这批训练结果，
算法就无法发现它的问题。如果假设空间中有许多假设，就必须考虑是否有至少一个
期望风险较高的假设也碰巧分对了全部训练邮件。因此，要判断需要多少训练邮件，既要计算
单个期望风险较高的假设出现这种情况的概率，也要考虑假设空间中有多少假设。

本章从有限假设空间的样本量计算出发，再处理阈值假设空间这样的无限集合。实数阈值虽然
有无限多个，在同一组有序输入上却只能形成少数几种预测。VC维刻画假设空间能在多少个
输入上给出全部可能的二元预测；一致收敛把这种表达能力与期望风险联系起来。
在此基础上，Rademacher复杂度与高斯复杂度进一步利用样本信息估计期望风险，样本压缩
和集成学习则从算法输出与组合方式建立保证。统计上存在学习器之后，还要检查能否用
可停机的程序实现它，以及训练和预测需要多少计算资源。

本章沿用第~\ref{sec:learning-problem}节的设定，把标签限制为两类，观测仍是输入与标签的配对，采用独立同分布抽样与零一损失。可实现与不可知PAC的要求见第~\ref{sec:pac-learning}节；这里用VC维、Rademacher复杂度和样本压缩，具体计算满足这些要求所需的样本量。通常把输出风险与假设空间$\mathcal H$内的最佳风险比较。集成与可计算学习中，若允许输出不属于$\mathcal H$的分类器，会明确说明这是非恰当学习。

\section{可实现的有限假设空间}
\label{sec:binary-finite-class}

回到邮件过滤的例子。假如算法只能从0.2、0.4、0.6和0.8这四个阈值中选择，
每个阈值对应一个分类假设，假设空间$\mathcal H$便只有四个成员。邮件分数仍可取
$[0,1]$中的任意值；有限的是可选假设的数量，而不是邮件分数的取值。
知道有哪些假设可选，就能逐个分析它们在样本上全部答对的可能性。

沿用可实现设定，假定$\mathcal H$中存在期望风险为零的假设$h^\star$。
随机抽取一封邮件，它出错的概率为零；因此它也几乎必然分对样本中的每一封邮件。
ERM选择经验风险最小的假设，至少可以选到一个经验风险为零的假设。我们称这样的
假设与样本\term{一致}（consistent），意思是它在这批样本上一次也没有分错。

但是，分对这批邮件的假设，面对新邮件时仍可能经常出错。假如一个假设分错随机新邮件
的概率为20\%，那么它分对一封邮件的概率就是80\%；独立抽取$n$封时，它碰巧全部
分对的概率是$0.8^n$。ERM只看得到这批样本，因而可能选中这样的假设。为了保证
ERM选出的假设期望风险较低，必须考虑假设空间中是否有任何一个期望风险较高的假设
碰巧与样本一致，而不能只计算其中某个假设的概率。

把20\%换成一般的误差容限$\varepsilon$。设$Y=\{0,1\}$、$S\sim\mathcal P^n$，
并取一个不随样本改变的假设$h\in\mathcal H$。如果它的期望风险
$R_{\mathcal P}(h)>\varepsilon$，随机抽取一封邮件，它被分对的概率就小于
$1-\varepsilon$。由于样本独立抽取，$h$与整批样本一致的概率满足
\begin{equation}
  \mathbb P\bigl(\hat R_S(h)=0\bigr)
  =\bigl(1-R_{\mathcal P}(h)\bigr)^n
  \leq(1-\varepsilon)^n\leq e^{-n\varepsilon}.
  \label{eq:bad-consistent-hypothesis}
\end{equation}
ERM会在看过样本后才决定选哪个假设，式~\eqref{eq:bad-consistent-hypothesis}不能直接
当作ERM选错的概率。不过，如果ERM选中了一个与样本一致、期望风险却超过
$\varepsilon$的假设，$\mathcal H$中就至少有一个这样的假设分对了整批样本。
对$\mathcal H$中的假设分别计算这一概率，再把结果相加，便得到“至少有一个”的
概率上界；即使几个假设同时分对样本，这个上界依然成立。由于$\mathcal H$只有
有限个假设，我们便能算出使这一概率足够小所需的样本量。

\begin{theorem}[有限假设空间的可实现PAC样本量界]
\label{thm:finite-realizable-pac}
设$\mathcal H$有限，$\mathcal P$对$\mathcal H$可实现。若算法输出任意满足
$\hat R_S(\hat h)=0$的$\hat h\in\mathcal H$，则
\begin{equation}
  \mathbb P\bigl(R_{\mathcal P}(\hat h)>\varepsilon\bigr)
  \leq |\mathcal H|e^{-n\varepsilon}.
  \label{eq:finite-realizable-failure}
\end{equation}
特别地，只要
\begin{equation}
  n\geq\frac{\log|\mathcal H|+\log(1/\delta)}{\varepsilon},
  \label{eq:finite-realizable-sample}
\end{equation}
算法就以至少$1-\delta$的概率输出期望风险不超过$\varepsilon$的假设。
\end{theorem}

\begin{proof}
若输出假设的期望风险大于$\varepsilon$，至少有一个期望风险大于$\varepsilon$的假设与样本一致。因此
\begin{align*}
 \mathbb P\bigl(R_{\mathcal P}(\hat h)>\varepsilon\bigr)
 &\leq \mathbb P\bigl(\exists h\in\mathcal H:
       R_{\mathcal P}(h)>\varepsilon,\ \hat R_S(h)=0\bigr)\\
 &\leq\sum_{h:R_{\mathcal P}(h)>\varepsilon}
       \mathbb P\bigl(\hat R_S(h)=0\bigr)
 \leq |\mathcal H|e^{-n\varepsilon}.
\end{align*}
令最后一项不超过$\delta$并整理$n$，得到式~\eqref{eq:finite-realizable-sample}。
\end{proof}

在四个阈值的例子中，$|\mathcal H|=4$。如果其中一个阈值对应的假设具有零期望风险，
我们希望ERM以至少99\%的概率选到期望风险不超过5\%的假设，就取
$\varepsilon=0.05$、$\delta=0.01$。式~\eqref{eq:finite-realizable-sample}给出
$\lceil(\log4+\log100)/0.05\rceil=120$。也就是说，独立同分布地抽取120封
训练邮件后，ERM选出的与样本一致的假设满足上述要求。这个数对所有满足可实现条件的
数据分布都适用，但不表示样本少于120封时一定会失败。

如果可选阈值从4个增加到8个，保持5\%的期望风险要求和99\%的成功概率不变，
公式给出的充分样本量只从120封增至
$\lceil(\log8+\log100)/0.05\rceil=134$封。一般地，假设数翻倍，式中尚未取整的
样本量只增加$\log2/\varepsilon$。因此，有限假设空间在可实现设定下是PAC可学习的：
对任意给定的$\varepsilon$和$\delta$，都能找到一个不依赖具体数据分布的充分样本量。
这里保证的是样本数量；能否高效地找到假设，还要另行分析。

这段推导依赖可实现设定：至少有一个假设的期望风险为零，因而几乎必然把训练样本全部分对。样本足够多时，期望风险
较高的假设也全部分对的机会就很小，ERM因而不容易选中它。标签有噪声时，所有假设
都可能在样本上出错，ERM只能比较它们各自错了多少。一个面对新邮件时经常出错的
假设，可能恰好在这批样本上错得很少。怎样判断样本上的比较在新邮件上仍然可靠？

\section{不可知的有限假设空间}
\label{sec:binary-uniform-convergence}

ERM看过样本后才选择经验风险最小的假设。某个假设即使面对新邮件时经常出错，也可能碰巧在这批样本上少犯错，从而被ERM选中。一种保证比较可靠的办法，是让同一批样本上所有假设的经验风险，都与各自的期望风险相差很小。这项要求称为一致收敛。
有限假设空间允许我们直接证明这项性质，并据此分析不可知设定下的ERM；无限假设空间
能否满足同样的要求，则留到第~\ref{sec:binary-infinite-class}节讨论。

设想我们先定好一个阈值，再抽取训练邮件。这个阈值对应的假设$h$不会随样本改变。
它分错随机新邮件的概率是期望风险$R_{\mathcal P}(h)$，在$n$封样本邮件中分错的
比例则是经验风险$\hat R_S(h)$。换一批邮件，它的经验风险可能变化。那它与期望风险相差
超过$\eta$的概率有多大？对每封样本邮件，$\ell_{01}(h(x_i),y_i)$取0或1，
表示$h$是否分错。由于样本邮件是独立抽取，这些随机变量相互独立，均值都是
$R_{\mathcal P}(h)$。Hoeffding不等式给出
\begin{equation}
 \mathbb P\left(\left|\hat R_S(h)-R_{\mathcal P}(h)\right|>\eta\right)
 \leq2e^{-2n\eta^2}.
 \label{eq:hoeffding-fixed-hypothesis}
\end{equation}
式~\eqref{eq:hoeffding-fixed-hypothesis}只控制抽样前定好的$h$，不能直接套在
随样本改变的ERM输出$A(S)$上。例如，用同一批带标签的邮件分别检验四个阈值，
我们可以算出这四个假设的经验风险，ERM选择其中经验风险最小的假设。若四个经验风险与各自的期望风险都相差不超过$\eta$，那么所选假设的期望风险至多比四者中的最低期望风险高$2\eta$。其中一份$\eta$来自对所选假设的估计，另一份来自对最佳比较假设的估计；后面的定理将逐步证明这一点。假设数量增加后，同样可以通过控制所有假设的估计偏差，保证选出的假设不会离最佳比较假设太远。
图~\ref{fig:learning-uniform-convergence}用竖线连接每个假设的经验风险与期望风险；
如果每条线都短，看过样本后才选出的假设也会受到保证。

\begin{figure}[htbp]
  \centering
  % 固定假设的集中与数据依赖选择之间的区别。
\begin{tikzpicture}[figure picture,foundation picture,x=1mm,y=1mm,font=\figurefont\footnotesize]
  \draw[foundation flow] (0,0) -- (94,0) node[right] {假设};
  \draw[foundation flow] (0,0) -- (0,48) node[above] {分类错误率};
  \draw[FigureGrid,line width=.5pt] (0,10) -- (92,10);
  \draw[FigureGrid,line width=.5pt] (0,25) -- (92,25);
  \draw[FigureGrid,line width=.5pt] (0,40) -- (92,40);

  \foreach \x/\pop/\emp in {10/31/29,25/22/19,40/27/23,55/18/12,70/24/20,85/20/17}{
    \draw[FigureRed,line width=1.1pt] (\x-3,\pop) -- (\x+3,\pop);
    \fill[FigureBlue] (\x,\emp) circle (1.3);
    \draw[FigureMuted,line width=.5pt] (\x,\emp) -- (\x,\pop);
  }
  \node[font=\figurefont\footnotesize,text=FigureRed,anchor=west] at (5,45) {横线：期望风险$R_{\mathcal P}(h)$};
  \node[font=\figurefont\footnotesize,text=FigureBlue,anchor=west] at (47,45) {圆点：经验风险$\hat R_S(h)$};
  \draw[FigureYellow,line width=1.1pt,rounded corners] (49,7) rectangle (61,34);
  \node[font=\figurefont\footnotesize,text=FigureInk,align=center] at (55,4)
    {ERM输出};
  \draw[foundation flow,FigureYellow] (55,7) -- (55,10.5);
\end{tikzpicture}%

  \caption{固定一个假设与同时考察全部假设的区别。即使某个假设的经验风险接近其期望风险，
  ERM仍可能选中其他估计偏差较大的假设。若全部竖线都短，框中由样本选出的假设也有保证。}
  \label{fig:learning-uniform-convergence}
\end{figure}

图中的每条竖线都短，意味着没有哪个假设的经验风险明显偏离期望风险。以四个阈值
为例，设可接受的风险差为$\eta$：只有四个阈值各自的经验风险与期望风险都相差不超过
$\eta$，这批邮件才满足“同时接近”。我们希望找到一个邮件数，使独立抽取这么多
带标签邮件时，出现这种情形的概率至少为$1-\delta$；这里$\delta$是允许的失败概率。
而且，这个邮件数不能随邮件与标签的分布改变。下面将这一要求写成正式定义。

\paragraph{可测性与可分性约定}

有限假设空间中的上确界只是有限个可测随机变量的最大值，不会产生额外的可测性问题。处理无限类时，本章对损失类
\[
 \mathcal F_{01}=\{(x,y)\mapsto\ell_{01}(h(x),y):h\in\mathcal H\}
\]
采用第~\ref{sec:learning-problem}节的可测性与可分性约定。压缩、重构和学习器输出涉及的数据依赖映射也须可测；可实现的零一损失样本可选择与样本一致的成员，其他存在性证明则使用该约定中的可测近似ERM。若这些条件不成立，本章结论不能无条件外推。

\begin{definition}[一致收敛]
\label{def:learning-uniform-convergence}
设$\mathcal H$是由可测函数$h:X\to\{0,1\}$组成的非空假设空间，采用零一损失并满足
第~\ref{sec:learning-problem}节的可测性与可分性约定。如果对
任意$\eta,\delta\in(0,1)$，都存在只依赖$\eta,\delta$和$\mathcal H$的正整数
$m_{\mathcal H}^{\mathrm{UC}}(\eta,\delta)$，使得对任意定义在$X\times\{0,1\}$上的
联合分布$\mathcal P$及任意$n\geq m_{\mathcal H}^{\mathrm{UC}}(\eta,\delta)$，独立同分布样本
$S\sim\mathcal P^n$满足
\begin{equation}
 \mathbb P_{S\sim\mathcal P^n}\left(
   \sup_{h\in\mathcal H}|\hat R_S(h)-R_{\mathcal P}(h)|\leq\eta
 \right)\geq1-\delta,
 \label{eq:learning-uniform-convergence-definition}
\end{equation}
就称$\mathcal H$满足\term{一致收敛}（uniform convergence）。
\end{definition}

式~\eqref{eq:learning-uniform-convergence-definition}中的上确界检查$\mathcal H$内全部假设的
估计偏差。它不超过$\eta$，表示这一批样本对每个假设都估得准。即使最大偏差
没有被某个假设取到，上确界仍然给出这些偏差的最小上界。概率是对整批样本
的随机抽取而言，不要求每一批样本都符合条件。这里$\mathcal P$同时规定输入$x$
出现的频率以及给定输入时标签$y$如何产生，并非只指输入的边缘分布$\mathcal P_X$。
不同分布可以改变各假设的期望风险，但不能改变保证所需的样本量
$m_{\mathcal H}^{\mathrm{UC}}(\eta,\delta)$。定义也不要求$\mathcal H$中存在零期望
风险假设，因此适用于不可知设定；它的分布无关要求与定义~\ref{def:general-pac-learnability}相同。

一致收敛保证的是估计准确，并不要求所有假设的期望风险都低：
期望风险高的假设只要经验风险也高，仍符合要求。

对有限假设空间，所有假设的经验风险同时接近期望风险，等价于没有任何假设的
风险差超过$\eta$。式~\eqref{eq:hoeffding-fixed-hypothesis}给出了每个固定假设
出现较大偏差的概率上界。有限个假设的这些概率相加，便能控制是否有任何一个
假设估计失准。

\begin{theorem}[有限假设空间的一致收敛样本量界]
\label{thm:finite-uniform-convergence}
设非空假设空间$\mathcal H$有限，$\mathcal P$是任意输入与标签的联合分布，
$S\sim\mathcal P^n$。对任意$\eta,\delta\in(0,1)$，只要样本量满足
\begin{equation}
 n\geq\frac{\log(2|\mathcal H|/\delta)}{2\eta^2},
 \label{eq:finite-uniform-sample}
\end{equation}
就有
\[
 \mathbb P_{S\sim\mathcal P^n}\left(
 \sup_{h\in\mathcal H}|\hat R_S(h)-R_{\mathcal P}(h)|\leq\eta
 \right)\geq1-\delta.
\]
同一概率界也可写为：对任意$n\geq1$和$\delta\in(0,1)$，以下不等式以至少$1-\delta$的概率成立：
\begin{equation}
 \sup_{h\in\mathcal H}|\hat R_S(h)-R_{\mathcal P}(h)|
 \leq\sqrt{\frac{\log(2|\mathcal H|/\delta)}{2n}}.
 \label{eq:finite-uniform-bound}
\end{equation}
因此，有限假设空间满足定义~\ref{def:learning-uniform-convergence}中的一致收敛。
\end{theorem}

\begin{proof}
对每个固定的$h\in\mathcal H$，式~\eqref{eq:hoeffding-fixed-hypothesis}给出
$|\hat R_S(h)-R_{\mathcal P}(h)|>\eta$的概率上界。只要其中一个假设发生这种偏差，
上确界就会超过$\eta$；把至多$|\mathcal H|$个概率相加，可得
\[
 \mathbb P\left(\sup_{h\in\mathcal H}|\hat R_S(h)-R_{\mathcal P}(h)|>\eta\right)
 \leq 2|\mathcal H|e^{-2n\eta^2}.
\]
若$n$满足式~\eqref{eq:finite-uniform-sample}，右侧至多为$\delta$。
反过来，令右侧等于$\delta$并解出$\eta$，得到式~\eqref{eq:finite-uniform-bound}。
两个界都对任意$\mathcal P$成立，所需样本量不依赖$\mathcal P$。
\end{proof}

一致收敛保证同一份样本能准确比较所有假设的风险。ERM输出$\hat h$的经验风险不高于
任意比较假设$h\in\mathcal H$的经验风险；若两者的经验风险都接近各自的期望风险，
便能把样本上的比较转成期望风险的比较。下面直接给出有限假设空间所需的样本量。

\begin{theorem}[有限假设空间的不可知PAC样本量界]
\label{thm:finite-agnostic-pac}
设非空假设空间$\mathcal H$有限，$\mathcal P$是任意输入与标签的联合分布，
$S\sim\mathcal P^n$，$\hat h$是$\mathcal H$上的精确ERM。对任意
$\varepsilon,\delta\in(0,1)$，只要
\begin{equation}
 n\geq\frac{2}{\varepsilon^2}
 \left(\log(2|\mathcal H|)+\log\frac1\delta\right),
 \label{eq:finite-agnostic-sample}
\end{equation}
就有
\[
 \mathbb P_{S\sim\mathcal P^n}\left(
 R_{\mathcal P}(\hat h)\leq\inf_{h\in\mathcal H}R_{\mathcal P}(h)+\varepsilon
 \right)\geq1-\delta.
\]
\end{theorem}

\begin{proof}
在定理~\ref{thm:finite-uniform-convergence}中取$\eta=\varepsilon/2$。
式~\eqref{eq:finite-agnostic-sample}保证
$\sup_{h\in\mathcal H}|\hat R_S(h)-R_{\mathcal P}(h)|\leq\varepsilon/2$
以至少$1-\delta$的概率成立。固定一份满足这一条件的样本，任取$h\in\mathcal H$。
ERM输出的$\hat h$也属于$\mathcal H$，所以两者的经验风险与期望风险都相差不超过
$\varepsilon/2$。由ERM的选择规则，$\hat R_S(\hat h)\leq\hat R_S(h)$，于是
\begin{align*}
 R_{\mathcal P}(\hat h)
 &\leq\hat R_S(\hat h)+\frac{\varepsilon}{2}\\
 &\leq\hat R_S(h)+\frac{\varepsilon}{2}\\
 &\leq R_{\mathcal P}(h)+\varepsilon.
\end{align*}
第一步和第三步分别使用输出假设与比较假设的风险估计，第二步使用ERM的选择规则。
这个不等式对每个$h\in\mathcal H$都成立，对右侧取下确界即得定理结论。
\end{proof}

这里的目标是让输出假设比假设空间内最低期望风险至多高$\varepsilon$。与可实现设定的
式~\eqref{eq:finite-realizable-sample}相比，充分界对误差容限的依赖从$1/\varepsilon$
变为$1/\varepsilon^2$。在可实现设定下，假设空间中有一个期望风险为零的假设，
它几乎必然分对全部训练样本；只需防止期望风险超过$\varepsilon$的假设也碰巧全部分对。
在不可知设定下，假设可能都在训练样本上犯错，ERM必须比较错误比例的大小。
因此，我们要求每个假设的经验风险都接近其期望风险，而式~\eqref{eq:hoeffding-fixed-hypothesis}
中这种估计偏差的概率随$n\eta^2$增大而下降，这便产生了上面的$1/\varepsilon^2$充分界。

同一个界也为定理~\ref{thm:srm-oracle-bound}中的复杂度惩罚提供了具体形式。
设$\mathcal H_1\subseteq\mathcal H_2\subseteq\cdots$是预先给定的一列有限假设空间，
其中$k$表示第几个空间。结构风险最小化会根据样本选择$k$和$h\in\mathcal H_k$。
如果只控制某个预先指定的$\mathcal H_k$中的估计偏差，算法仍可能选中另一个空间：
那个空间中某些假设的经验风险恰好低于期望风险，因而看起来特别好。因此，我们需要
控制同一份样本对每个$\mathcal H_k$中所有假设的估计偏差。

具体地，在抽样前选定$\delta_k>0$，使$\sum_k\delta_k\leq\delta$。第$k$个空间
允许的估计偏差记为
\[
 p_k=\sqrt{\frac{\log(2|\mathcal H_k|/\delta_k)}{2n}}.
\]
式~\eqref{eq:finite-uniform-bound}
保证：对单独一个$\mathcal H_k$，其中某个假设满足
$|\hat R_S(h)-R_{\mathcal P}(h)|>p_k$的概率至多为$\delta_k$。把各空间的这些
失败概率相加，可知至少以$1-\delta$的概率，同一份样本满足
\[
 |\hat R_S(h)-R_{\mathcal P}(h)|\leq p_k
 \qquad(k\geq1,\ h\in\mathcal H_k).
\]
这个不等式对每个空间中的每个假设同时成立。无论算法后来选中哪个$(k,h)$，其期望风险都不超过
经验风险加上对应的$p_k$。因此，结构风险最小化可以选择使
\begin{equation}
 \hat R_S(h)+p_k
 \label{eq:srm-finite-penalty}
\end{equation}
最小的$(k,h)$。空间越大，或者分配给它的$\delta_k$越小，$p_k$就越大。例如，
取$\delta_k=\delta/[k(k+1)]$时，因$1/[k(k+1)]=1/k-1/(k+1)$，所有$\delta_k$
之和恰为$\delta$；这是一种无需额外选择规则的分配方法。

另一种分配方法可以把事先选定的空间描述长度写进惩罚项。这里要描述的只是序列中选择了
哪个$\mathcal H_k$，也就是它的序号$k$。我们可以预先给每个序号一个二进制编码，并要求没有
任何编码是另一编码的开头，这称为前缀编码。例如，$0$、$10$、$11$可以分别标识
三个空间。任取有限组编码，并考察比它们都长的二进制串：以某个$L_k$位编码开头的串
占同长度串的$2^{-L_k}$，不同前缀编码对应的串又互不重叠。因此，这组编码的
$2^{-L_k}$之和不超过1；任意有限组都如此，全部编码也满足
$\sum_k2^{-L_k}\leq1$，所以
取$\delta_k=\delta2^{-L_k}$仍满足总失败概率的要求。此时
$\log(1/\delta_k)=\log(1/\delta)+L_k\log2$：较长的编码通过较小的$\delta_k$
增大对应的$p_k$。前缀编码的作用是提供一种有效的失败概率分配，并使选择空间所需的
编码长度成为惩罚的一部分；前面的PAC保证本身不要求使用这种编码。

\section{无限假设空间}
\label{sec:binary-infinite-class}

有限假设空间中，我们分别计算每个假设的经验风险与期望风险相差超过容限的概率，再把这些概率
相加，得到式~\eqref{eq:finite-uniform-sample}的一致收敛样本量界。如果假设空间无限，
直接按假设个数相加便无法得到这样的上界。要判断无限假设空间能否一致收敛，需要
换一种方式衡量它的预测能力：考察同一组有限输入上，空间内的假设能给出多少种不同答案。

\subsection{VC维与生长函数的学习论接口}
\label{sec:binary-shattering-vc}

打散、VC维和生长函数的正式定义、有限模式例与示意图已经放在
\BookChapterRef{chap:combinatorics}{第一卷“数学准备”的“计数与计算”章}。
本章沿用其中的\term{限制}（restriction）类$\mathcal H|_C$与
\term{生长函数}（growth function）记号
\begin{equation}
 \Pi_{\mathcal H}(m)
 =\max_{x_1,\ldots,x_m\in X}
   \left|\{(h(x_1),\ldots,h(x_m)):h\in\mathcal H\}\right|.
 \label{eq:learning-growth-function}
\end{equation}
即$m$个输入上最多能实现的不同二元预测数；输入可重复的约定使有限输入空间上也能对任意
$m$使用该记号。第一卷已经算出阈值类、单区间类和平面仿射半空间的VC维分别为
$1,2,3$，以及前两类的精确生长函数。本章接着计算：这些有限的VC维，能为经验风险与期望风险的差提供什么保证，又需要多少训练样本？为此，先用下面的组合上界控制有限输入上的预测数。

\subsection{Sauer--Shelah引理}
\label{sec:binary-growth-sauer}

VC维只告诉我们，假设空间最多能在多少个输入上实现全部标签组合；风险分析却要知道更大的样本上还能出现多少种预测。Sauer--Shelah引理把前一个限制转成后一个计数上界。
\begin{lemma}[Sauer--Shelah引理]
\label{lem:learning-sauer-shelah}
若$1\leq d=\operatorname{VC}(\mathcal H)<\infty$，则对$m\geq d$，
\begin{equation}
 \Pi_{\mathcal H}(m)
 \leq\sum_{j=0}^{d}\binom{m}{j}
 \leq\left(\frac{em}{d}\right)^d.
 \label{eq:learning-sauer-shelah}
\end{equation}
\end{lemma}
\begin{proof}
令$G(m,d)$表示VC维至多为$d$的二分类假设空间在$m$个输入上能够实现的最大预测数。
对$m+d$作强归纳。边界情形为$G(0,d)=1$；当$d=0$时，任意一个输入都不能同时取两个
标签，所以$G(m,0)=1$。

现在固定$m\geq1,d\geq1$以及$m$个互异输入，并令
$V\subseteq\{0,1\}^{m}$为相应的预测向量集合。记$U$为删去最后一个坐标所得的前缀集合，
再记
\[
 U'=\{u\in\{0,1\}^{m-1}:(u,0)\in V,\ (u,1)\in V\}.
\]
每个$u\in U$至少贡献一个扩展，而恰有$u\in U'$再多贡献一个扩展，因此
$|V|=|U|+|U'|$。集合$U$的VC维至多为$d$；$U'$的VC维至多为$d-1$，否则若$U'$打散
前$m-1$个坐标中的$d$个，再加入最后一个坐标，$V$便会打散$d+1$个坐标。强归纳假设于是给出
\[
 G(m,d)\leq G(m-1,d)+G(m-1,d-1)
 \leq\sum_{j=0}^{d}\binom{m-1}{j}
   +\sum_{j=0}^{d-1}\binom{m-1}{j}.
\]
使用Pascal恒等式并约定$j>m$时$\binom mj=0$，得到
$G(m,d)\leq\sum_{j=0}^{d}\binom mj$，从而得到第一道不等式。

第二道不等式只把组合数写成便于代入样本界的形式。令$a=d/m\leq1$，则对$j\leq d$
有$a^j\geq a^d$，于是
\[
 a^d\sum_{j=0}^d\binom mj
 \leq\sum_{j=0}^d\binom mj a^j
 \leq(1+a)^m\leq e^{ma}=e^d.
\]
两边除以$a^d$即得结论。若$d=0$，非空假设空间在每个点上的标签都唯一，生长函数恒为1，
无需使用含有$d$作为分母的表达式。
\end{proof}

当$d$固定时，引理右侧随$m$至多按$m^d$的量级增长，而全部二元预测的数目是$2^m$。
这说明有限VC维足以阻止生长函数一直按指数速度增长；实际增长仍可能比引理的上界更慢，
如VC维为3的平面仿射半空间就只按$m^2$的量级增长。

\BookChapterRef{chap:combinatorics}{第一卷“数学准备”的“计数与计算”章}
还在同一对数坐标中比较了全部$2^m$种标签组合与有限VC维带来的多项式计数。
曲线间不断增大的差距表示可实现模式在全部二元组合中所占比例持续缩小，
并不要求假设空间本身只有有限多个函数。

\subsection{二分类可学习性基本定理}
\label{sec:binary-learnability-characterization}

我们希望建立经验风险与期望风险之间的关系：算法在样本上找到经验风险较小的假设后，
需要保证这个假设的期望风险也足够小。前面的一致收敛分析表明，只要同一份样本上
所有假设的经验风险与各自的期望风险都足够接近，ERM就能得到这样的保证。
对于无限假设空间，接下来要说明生长函数如何帮助我们控制这些差值。

期望风险取决于未知的数据分布，无法直接计算。为分析它与经验风险的差，我们在证明中
引入另一份样本$S'=(z_1',\ldots,z_n')$：它与原样本$S=(z_1,\ldots,z_n)$相互独立，
每个观测都来自同一个输入与标签的联合分布$\mathcal P$，其中$z_i=(x_i,y_i)$。
这份\term{副样本}（ghost sample）只用于理论分析，学习算法不需要实际取得它。

给定原样本$S$并选定一个假设$h$后，可以在副样本上重新计算它的经验风险
$\hat R_{S'}(h)$。由于副样本没有参与选择$h$，它能独立评估这个假设的预测错误率；
对副样本的所有可能抽样结果取平均，便得到$R_{\mathcal P}(h)$。
如果$\hat R_S(h)$与$R_{\mathcal P}(h)$相差很大，而$\hat R_{S'}(h)$接近
$R_{\mathcal P}(h)$，那么两份样本上的经验风险也会相差较大。因此，通过控制
两份样本上经验风险的差，可以间接控制原样本的经验风险与期望风险的差。
这里并不要求副样本已经能同时准确评估所有假设；证明只需对原样本确定后选出的
一个假设使用独立抽样的估计，再将两种偏差事件的概率联系起来。

两份样本的比较只涉及$2n$个观测，这就使生长函数能够发挥作用。在这些输入上，
给出完全相同预测的假设，在两份样本上的经验风险也完全相同。因此，即使假设有
无限多个，需要区分的预测结果也至多有$\Pi_{\mathcal H}(2n)$种。
接着，将每对观测$(z_i,z_i')$以相同概率保持原位或相互交换，各对独立进行。
两份样本独立同分布，所以这种交换不改变样本的联合分布；对一种固定的预测结果，
交换只会随机改变每对观测对经验风险差值的贡献的正负号。由此可以用集中不等式
控制这一差值较大的概率，再把至多$\Pi_{\mathcal H}(2n)$种预测结果对应的概率相加，
得到存在任意一个假设使差值过大的概率上界。引入副样本并进行这种随机交换，
构成了\term{对称化}（symmetrization）分析。

下面先用副样本把未知期望换成两个样本均值的比较，再用随机交换控制每一种预测的偏差，最后用生长函数统计需要同时控制多少种预测。

\begin{theorem}[VC一致收敛界]
\label{thm:vc-uniform-convergence}
设$\mathcal H$满足本章的可测性约定，$S\sim\mathcal P^n$，并令
$d=\operatorname{VC}(\mathcal H)$。对任意$\eta>0$，只要$n\eta^2\geq2$，就有
\begin{equation}
 \mathbb P\left\{
   \sup_{h\in\mathcal H}|R_{\mathcal P}(h)-\hat R_S(h)|>\eta
 \right\}
 \leq4\Pi_{\mathcal H}(2n)e^{-n\eta^2/8}.
 \label{eq:vc-symmetrization-probability}
\end{equation}
因而当$1\leq d\leq2n$时，以至少$1-\delta$的概率，
\begin{equation}
 \sup_{h\in\mathcal H}|R_{\mathcal P}(h)-\hat R_S(h)|
 \leq
 \sqrt{\frac8n\left[
   d\log\frac{2en}{d}+\log\frac4\delta
 \right]}.
 \label{eq:vc-uniform-bound}
\end{equation}
若右侧大于1，则可用零一损失给出的平凡上界1替代。
\end{theorem}

\begin{proof}
取独立副样本$S'\sim\mathcal P^n$。当$n\eta^2\geq2$时，固定$S$并选择使原偏差超过
$\eta$的假设。零一损失的样本均值方差不超过$1/(4n)$，Chebyshev不等式因而把副样本风险偏离期望风险至少$\eta/2$的条件概率控制在$1/(n\eta^2)\leq1/2$。
于是，副样本风险以至少$1/2$的条件概率落在期望风险的$\eta/2$以内，标准对称化步骤给出
\[
 \mathbb P\left\{\sup_h|R_{\mathcal P}(h)-\hat R_S(h)|>\eta\right\}
 \leq
 2\mathbb P\left\{\sup_h|\hat R_{S'}(h)-\hat R_S(h)|>\frac\eta2\right\}.
\]
条件于合并后的$2n$个输入，假设只产生至多$\Pi_{\mathcal H}(2n)$种不同预测向量。
对每种固定预测向量，记$g_h(z)=\ell_{01}(h(x),y)$。独立交换$(z_i,z_i')$，等价于给第$i$对的损失差乘以独立Rademacher符号$\sigma_i$。给定两份原样本后，损失差已确定，随机性只来自这些符号；每项除以$n$后的取值区间长度至多为$2/n$。Hoeffding不等式于是给出
\[
 \mathbb P_\sigma\left\{
   \left|\frac1n\sum_{i=1}^n\sigma_i
     \bigl(g_h(z_i')-g_h(z_i)\bigr)\right|>\frac\eta2
   \,\middle|\,S,S'
 \right\}
 \leq2e^{-n\eta^2/8}.
\]
对至多$\Pi_{\mathcal H}(2n)$种向量使用并集界，再乘上对称化的因子2，得到
式~\eqref{eq:vc-symmetrization-probability}。最后由Sauer--Shelah引理代入
$\Pi_{\mathcal H}(2n)\leq(2en/d)^d$，并令右侧不超过$\delta$，得到
式~\eqref{eq:vc-uniform-bound}。
\end{proof}

固定有限的$d$和失败概率容限$\delta$后，式~\eqref{eq:vc-uniform-bound}右侧随$n$增大
趋于零。因此，只要样本量足够大，就能以至少$1-\delta$的概率，使所有假设的经验风险
与各自的期望风险之差都不超过指定容限。同一个样本量要求适用于所有允许的联合分布
$\mathcal P$，这就得到了一致收敛。

由前面ERM的分析，一致收敛进一步保证不可知PAC可学习。反过来，若VC维无限，
没有免费的午餐定理会排除分布无关的PAC学习保证。由此，VC维、一致收敛与PAC
可学习性之间形成了下面的等价关系。式~\eqref{eq:vc-uniform-bound}足以证明
这种等价关系，但其中的对数因子不一定是最优样本复杂度所必需的。

\begin{theorem}[二分类学习基本定理]
\label{thm:binary-fundamental-theorem}
在二分类、零一损失、独立同分布抽样、分布无关且满足共享可测性约定的设定中，
对非空假设空间$\mathcal H$，下列命题等价：
\begin{enumerate}
 \item $\mathcal H$是不可知PAC可学习的；
 \item $\mathcal H$是可实现PAC可学习的；
 \item $\mathcal H$满足一致收敛；
 \item $\operatorname{VC}(\mathcal H)<\infty$。
\end{enumerate}
\end{theorem}
\begin{proof}
若VC维有限，Sauer--Shelah引理限制合并样本上的不同预测数；对称化把这种计数变成
式~\eqref{eq:vc-uniform-bound}的一致收敛界。在一致收敛事件上，ERM输出与任意比较
假设的经验风险之差不大于零，而两者的经验风险与期望风险各相差至多$\eta$，因此输出
假设的期望风险至多比$\inf_{h\in\mathcal H}R_{\mathcal P}(h)$高$2\eta$。令$\eta=\varepsilon/2$，
便得到不可知PAC可学习；不可知保证覆盖可实现分布，因而也推出可实现PAC可学习。

反过来，若VC维无限，对任意样本量$n$都能找到至少$2n$个可被$\mathcal H$打散的输入。
对任意给定的算法，第~\ref{sec:no-free-lunch}节的没有免费的午餐定理都能在这些输入上构造一个可实现分布，
使该算法只用$n$个观测时，以不低于某个正常数的概率输出期望风险超过固定误差阈值的假设。这与可实现PAC要求的
统一样本量矛盾。因此，可实现PAC可学习必然要求VC维有限。两条方向合起来得到四个命题
的等价性。
\end{proof}

这个定理把第~\ref{chap:statistical-learning-framework}章的PAC定义变成一个可检查的容量条件。
它说的是是否存在分布无关的学习保证；样本量是否达到最优阶，是另一项更细的要求。

前面的界给出“这些样本已经足够”的保证，却还未说明是否可以更少。最优样本复杂度把两方面放在一起：某个学习器在所有允许分布上需要多少样本，以及任何学习器在最难分布上至少需要多少样本。上下界达到同一量级，才能称为最优阶。
\begin{theorem}[二分类的最优样本复杂度阶]
\label{thm:binary-optimal-sample-complexity}
设$d=\operatorname{VC}(\mathcal H)\geq2$，并把$\varepsilon,\delta$限制在小于某个普适
正常数的范围内。在分布无关的极小极大意义下，允许非恰当学习时，可实现PAC学习的
最优样本复杂度阶为
\[
 \Theta\left(\frac{d+\log(1/\delta)}{\varepsilon}\right),
\]
不可知PAC学习的最优样本复杂度阶为
\[
 \Theta\left(\frac{d+\log(1/\delta)}{\varepsilon^2}\right).
\]
这里的$\Theta$略去与$\mathcal H,\varepsilon,\delta$无关的普适常数。
\end{theorem}

上述结果包含匹配的上界与下界，不能由定理~\ref{thm:vc-uniform-convergence}单独推出。
经典生长函数证明给出的可实现ERM上界会多出$\log(1/\varepsilon)$；更精细的非恰当
学习器可以去掉这项\scite{hanneke2016optimal,shalevshwartz2014}。因此，最优阶不表示
任意ERM都达到同样的样本量，也不表示要求输出属于$\mathcal H$时总能达到该阶。
维度为0或1的边界情形需要另行分析。这个统计结论也没有说明学习器是否容易计算；
计算问题将在本章末尾单独讨论。

\section{估计期望风险}
\label{sec:rademacher-complexity}

上一节利用VC维回答了二分类假设空间是否具有分布无关PAC可学习性。
接下来，我们关心一个更具体的问题：算法已经根据样本选出了一个假设，
它在未来数据上的期望风险最多有多大？经验风险可以由已有样本计算，
而期望风险依赖未知的数据分布，因此还需要估计两者之间可能存在的差距。

VC维给出的一致收敛结论也能用于估计这个差距，但它依据整个输入空间的打散能力，
没有直接利用当前样本的位置以及实值预测和损失的大小。
本节介绍Rademacher复杂度和高斯复杂度，用这些信息进一步分析期望风险。
所得结论具有共同形式：在可计算的经验风险上，加上复杂度项和抽样误差项，
得到一个以高概率成立的期望风险上界。这里的“估计”指给出这样的上界，
并不意味着能够从有限样本精确算出未知的期望风险。

\subsection{Rademacher复杂度}

Rademacher复杂度的定义、随机符号图示、凸包性质和线性类计算已经在
\BookSectionRef{sec:probability-uniform-errors-symmetry}{第一卷“随机变量”章的“候选族的共同误差”一节}
建立。它考察一批随机正负号产生后，候选函数能把样本上的函数值与这些符号配合得多好：若能从候选中挑出让许多乘积为正的函数，随机和的上确界就会变大。本章采用不带绝对值的单侧版本，因为风险上界先只需要控制“经验风险低估期望风险”的方向。对$S=(z_1,\ldots,z_n)$及非空实值函数类$\mathcal G$，记
\begin{equation}
 \widehat{\mathfrak R}_S(\mathcal G)
 =\mathbb E_\sigma\left[\sup_{f\in\mathcal G}
       \frac1n\sum_{i=1}^n\sigma_i f(z_i)\right].
 \label{eq:learning-empirical-rademacher}
\end{equation}
其中$\sigma_i$独立且等概率取$-1,1$；再记
\[
 \mathfrak R_n(\mathcal G)=\mathbb E_{S\sim\mathcal P^n}
       \widehat{\mathfrak R}_S(\mathcal G).
\]
第一卷使用带绝对值的对称版本；当$\mathcal G$关于取负封闭时两者相同，一般情况下
带绝对值版本不小于当前单侧版本。复用收缩或比较定理时必须按所用版本核对常数。

\subsection{线性预测的复杂度}

多维线性模型也能用同样的方法计算，并得到适用于任意样本量的上界。
考虑线性得分$f_{\bm w}(\bm x)=\bm w^\top\bm x$。若不限制参数大小，
当$\sum_i\sigma_i\bm x_i\neq0$时，可以不断增大该方向上的参数，
使$\sum_i\sigma_i f_{\bm w}(\bm x_i)$任意大。因此，需要对允许的参数作出限制。
设$B>0$，要求$\|\bm w\|_2\leq B$，即参数向量的长度不超过$B$。
这里先计算预测得分的复杂度，再将结果用于损失函数；为此，令
$\mathcal F=\{\bm{x}\mapsto \bm{w}^\top \bm{x}:\lVert\bm{w}\rVert_2\leq B\}$，
并假定样本中的每个输入满足$\lVert\bm{x}_i\rVert_2\leq r$。
记输入样本为$S_X=(\bm x_1,\ldots,\bm x_n)$。将线性得分代入定义，可得
\begin{align}
 \widehat{\mathfrak R}_{S_X}(\mathcal F)
 &=\frac1n\mathbb E_\sigma\sup_{\lVert\bm{w}\rVert_2\leq B}
   \bm{w}^\top\sum_i\sigma_i \bm{x}_i\nonumber\\
 &\leq\frac{B}{n}\sqrt{\sum_i\lVert\bm{x}_i\rVert_2^2}
 \leq\frac{Br}{\sqrt n}.
 \label{eq:linear-rademacher-bound}
\end{align}
这个界没有显式出现参数维数。固定随机符号后，令参数沿着
$\sum_i\sigma_i\bm x_i$的方向取长度$B$，最大内积就是该向量范数的$B$倍。
期望范数不超过均方范数的平方根，而独立符号使不同样本间的交叉项平均为零，所以只剩
$\sum_i\lVert\bm x_i\rVert_2^2$。这里的维数无关保证依赖参数半径$B$与输入半径$r$均受控，
不能在没有这些条件时沿用。

\subsection{由Rademacher复杂度估计期望风险}

前面已经说明如何定义和计算复杂度，现在把它与期望风险联系起来。
计算线性预测的复杂度时，函数值是预测得分；估计期望风险时，需要分析的函数值则是损失。
设损失函数$\ell$取值于$[0,1]$，对每个假设$h$定义
$g_h(z)=\ell(h(x),y)$，其中$z=(x,y)$，并令$\mathcal G=\{g_h:h\in\mathcal H\}$。
于是$g_h$在样本上的平均值是$\hat R_S(h)$，它在分布$\mathcal P$下的期望是
$R_{\mathcal P}(h)$。控制这两个平均值之差，就能由经验风险得到期望风险上界。

\begin{theorem}[Rademacher复杂度给出的期望风险上界]
\label{thm:rademacher-risk-bound}
设$\mathcal G$是抽样前确定的、取值于$[0,1]$的非空函数空间，并满足本章上述可测性
约定。若$S\sim\mathcal P^n$，则对任意$\delta\in(0,1)$，
以至少$1-\delta$的概率，对所有$g\in\mathcal G$同时有
\begin{equation}
 \mathbb E_{z\sim\mathcal P}g(z)
 \leq\frac1n\sum_i g(z_i)+2\mathfrak R_n(\mathcal G)
 +\sqrt{\frac{\log(1/\delta)}{2n}}.
 \label{eq:rademacher-risk-bound}
\end{equation}
另有一个由样本确定的上界：以至少$1-\delta$的概率，对所有$g\in\mathcal G$同时有
\begin{equation}
 \mathbb E_{z\sim\mathcal P}g(z)
 \leq\frac1n\sum_i g(z_i)+2\widehat{\mathfrak R}_S(\mathcal G)
 +3\sqrt{\frac{\log(2/\delta)}{2n}}.
 \label{eq:empirical-rademacher-risk-bound}
\end{equation}
\end{theorem}

\begin{proof}
记$\Phi(S)=\sup_{g\in\mathcal G}(\mathbb E g-\frac1n\sum_i g(z_i))$。
它表示在这份样本上，经验均值对真实期望的最大低估程度。引入独立副样本$S'$后，
先对副样本取平均再选择函数，不会大于先选择函数再取平均，因此
\[
 \mathbb E_S\Phi(S)
 \leq\mathbb E_{S,S'}\sup_{g\in\mathcal G}
 \frac1n\sum_i\bigl(g(z_i')-g(z_i)\bigr).
\]
对每对观测随机交换，可在差值前加入独立符号$\sigma_i$而不改变上述期望。
将两份样本对应的和分别取上确界，并利用$\sigma_i$与$-\sigma_i$同分布，得到
\begin{equation}
 \mathbb E_S\Phi(S)\leq2\mathfrak R_n(\mathcal G).
 \label{eq:rademacher-symmetrization}
\end{equation}

更换$S$中的一个观测，$\Phi(S)$最多改变$1/n$。有界差分不等式于是保证，
以至少$1-\delta$的概率，$\Phi(S)$不超过其期望加上
$\sqrt{\log(1/\delta)/(2n)}$，这就证明了第一个上界。
经验Rademacher复杂度在更换一个观测时也最多改变$1/n$，因而以至少$1-\delta/2$的概率，
\[
 \mathfrak R_n(\mathcal G)
 \leq\widehat{\mathfrak R}_S(\mathcal G)
 +\sqrt{\frac{\log(2/\delta)}{2n}}.
\]
将第一个上界的失败概率也设为$\delta/2$，这两个不等式同时成立的概率至少为
$1-\delta$。代入即得第二个上界。
\end{proof}

令$g=g_h$，定理左侧就是假设$h$的期望风险，右侧第一项就是它的经验风险。
其余两项说明，在多大程度上可以用经验风险代替期望风险。由于结论对所有假设同时成立，
即使$h$由算法观察样本后选出，也能使用这个上界。第二个上界还允许根据当前样本
计算或估计复杂度；两个假设空间即使具有相同VC维，也可能因在这批样本上的函数值不同，
而得到不同的期望风险上界。这类依赖数据的分析是Rademacher复杂度的重要用途
\scite{bartlettmendelson2002}。

线性预测的得分不是零一损失。若某个损失函数对预测值满足$L$-Lipschitz条件，
即预测值改变$a$时损失至多改变$L|a|$，\term{收缩引理}（contraction lemma）
可以把预测函数空间的复杂度界转给相应的损失函数空间。第一卷已经在
\BookSectionRef{sec:probability-uniform-errors-symmetry}{“随机变量”章的“候选族的共同误差”一节}
证明单侧与带绝对值版本；本章采用单侧复杂度，并在损失不满足$\phi(0)=0$时先减去
常数$\phi(0)$，该平移不改变随机和的期望。这个步骤需要核对损失的变化率，不能直接把
式~\eqref{eq:linear-rademacher-bound}当成零一损失下的期望风险界。

这个结果可以具体用于分类。将标签写成$y\in\{-1,1\}$，令$f(\bm x)=\bm w^\top\bm x$。
乘积$yf(\bm x)$称为带符号的分类间隔。这里的间隔位于得分尺度：它为正表示预测方向
正确，数值越大，分类得分距离决策阈值0越远；若要换成输入空间中的几何间隔，还需
除以参数向量的范数。给定$\rho>0$，定义间隔损失
\[
 \ell_\rho(f(\bm x),y)=
 \begin{cases}
  1,&yf(\bm x)\leq0,\\
  1-yf(\bm x)/\rho,&0<yf(\bm x)<\rho,\\
  0,&yf(\bm x)\geq\rho.
 \end{cases}
\]
预测错误时该损失为1；预测正确但间隔不足$\rho$时仍有一定惩罚。
它取值于$[0,1]$、关于得分是$1/\rho$-Lipschitz的，并且不小于零一损失。
若$\|\bm w\|_2\leq B$且$\|\bm x\|_2\leq r$几乎处处成立，
收缩引理和式~\eqref{eq:linear-rademacher-bound}给出间隔损失空间的复杂度上界
$Br/(\rho\sqrt n)$。于是，以至少$1-\delta$的概率，所有这些线性预测函数均满足
\begin{equation}
 R_{\mathcal P}^{01}(\operatorname{sign}f)
 \leq\frac1n\sum_i\ell_\rho(f(\bm x_i),y_i)
 +\frac{2Br}{\rho\sqrt n}
 +\sqrt{\frac{\log(1/\delta)}{2n}}.
 \label{eq:rademacher-margin-risk}
\end{equation}
这里$R_{\mathcal P}^{01}$明确表示零一损失下的期望风险，得分为0时采用固定分类规则。
若样本上的间隔损失为零，要使右侧不超过$\varepsilon$，所需样本量的充分上界为
$O(((Br/\rho)^2+\log(1/\delta))/\varepsilon^2)$。
这个保证由范数与间隔的比值控制，没有显式依赖输入维数；在高维且间隔较大的情形，
它可能比仅依据半空间VC维得到的保证更有用。

这里增加了范数受限、输入有界和间隔损失较小的条件，因此并没有推翻二分类学习基本定理。
单纯缩小参数范数而保持分类符号不变，通常也会缩小间隔，不能据此无条件改善上界。
Rademacher复杂度带来的改进，是能够利用这些具体条件，并把分析扩展到实值损失；
它不保证对每个问题都给出比VC维更小的样本量，也不能使无限VC维的二分类假设空间
在所有分布上变得PAC可学习。


要进一步得到ERM的可学习性保证，需要控制两个方向的偏差。将第一个上界分别用于
$\mathcal G$和$1-\mathcal G=\{1-g:g\in\mathcal G\}$。两个函数类都取值于
$[0,1]$；平移不改变期望均值与经验均值之差，而随机符号的对称性给出
$\mathfrak R_n(1-\mathcal G)=\mathfrak R_n(\mathcal G)$。分别分配失败概率
$\delta/2$，可知以至少$1-\delta$的概率，所有假设均满足
\[
 |R_{\mathcal P}(h)-\hat R_S(h)|
 \leq2\mathfrak R_n(\mathcal G)+\sqrt{\frac{\log(2/\delta)}{2n}}.
\]
若$\hat h$是ERM输出，按照前文比较两次经验风险与期望风险的方法，就得到
\begin{equation}
 R_{\mathcal P}(\hat h)-\inf_{h\in\mathcal H}R_{\mathcal P}(h)
 \leq4\mathfrak R_n(\mathcal G)
 +2\sqrt{\frac{\log(2/\delta)}{2n}}.
 \label{eq:rademacher-excess-risk}
\end{equation}
因此，若对某个分布有$\mathfrak R_n(\mathcal G)\to0$，ERM的期望风险就以高概率
接近假设空间所能达到的最小期望风险。若进一步存在对所有允许分布都成立的上界
$\mathfrak R_n(\mathcal G)\leq a_n$，且$a_n\to0$，便能选出不依赖具体分布的
样本量，得到不可知PAC保证。仅在某个分布下复杂度很小，并不足以推出分布无关PAC可学习。

\subsection{高斯复杂度}

经验高斯复杂度的正式定义及其与Rademacher复杂度的比较已经在
\BookSectionRef{sec:probability-uniform-errors-symmetry}{第一卷“随机变量”章的“候选族的共同误差”一节}
建立。本章沿用其中的独立标准高斯系数$\xi_i$，并将单侧
\term{经验高斯复杂度}（empirical Gaussian complexity）记为
\[
 \widehat{\mathfrak G}_S(\mathcal G)
 =\mathbb E_\xi\left[\sup_{g\in\mathcal G}\frac1n\sum_i \xi_i g(z_i)\right].
\]
高斯系数的幅度可变，Rademacher符号的幅度恒为1，因此一般不能断言两种复杂度只差
普适常数。对当前单侧版本，写成$\xi_i=\sigma_i|\xi_i|$并利用上确界的凸性可得
$\widehat{\mathfrak G}_S(\mathcal G)\geq\sqrt{2/\pi}\,\widehat{\mathfrak R}_S(\mathcal G)$。
因此，若已得到$\widehat{\mathfrak G}_S(\mathcal G)\leq b(S)$，便有
$\widehat{\mathfrak R}_S(\mathcal G)\leq\sqrt{\pi/2}\,b(S)$，可代入
式~\eqref{eq:empirical-rademacher-risk-bound}来控制期望风险。

这给出了直接使用高斯复杂度的期望风险上界。

\begin{theorem}[高斯复杂度给出的期望风险上界]
\label{thm:gaussian-risk-bound}
设$\mathcal G=\{g_h:h\in\mathcal H\}$为抽样前确定的、取值于$[0,1]$的损失函数空间，
并满足定理~\ref{thm:rademacher-risk-bound}的可测性条件。
若$S\sim\mathcal P^n$，则对任意$\delta\in(0,1)$，以至少$1-\delta$的概率，
所有$h\in\mathcal H$同时满足
\begin{equation}
 R_{\mathcal P}(h)\leq\hat R_S(h)
 +\sqrt{2\pi}\,\widehat{\mathfrak G}_S(\mathcal G)
 +3\sqrt{\frac{\log(2/\delta)}{2n}}.
 \label{eq:gaussian-risk-bound}
\end{equation}
\end{theorem}

\begin{proof}
在式~\eqref{eq:empirical-rademacher-risk-bound}中取$g=g_h$，并代入
$2\widehat{\mathfrak R}_S(\mathcal G)\leq\sqrt{2\pi}\,\widehat{\mathfrak G}_S(\mathcal G)$，
即得结论。
\end{proof}

实际使用时，先计算所选假设的经验风险，再为损失函数空间的经验高斯复杂度求一个上界，
将两者代入式~\eqref{eq:gaussian-risk-bound}即可。若已有复杂度上界$b(S)$，
也可以直接用$b(S)$替换其中的经验高斯复杂度。这个结论对所有假设同时成立，
因此同样适用于观察样本后选出的假设。

高斯复杂度在这里提供了估计复杂度项的另一种办法，并不保证得到更小的期望风险上界。
反过来用Rademacher复杂度控制高斯复杂度时，所需的倍数可能随$n$增长。例如限制向量仅为$\{\pm\bm e_i:1\leq i\leq n\}$时，
Rademacher复杂度为$1/n$，高斯复杂度为$\mathbb E\max_i|\xi_i|/n$。两种工具联系紧密，
但互换两类复杂度得到的期望风险界时仍须检查具体比较条件。

\section{样本压缩}
\label{sec:sample-compression}

前两节分别用VC维和Rademacher复杂度分析假设空间，进而控制经验风险与期望风险的关系。
还有一种办法从学习算法的输出入手：如果只保留少量带标签观测，就能重新构造出一个
与全部样本一致的假设，那么这些被保留的观测能否为它的期望风险提供保证？
\term{样本压缩}（sample compression）研究的正是这种从少量观测重构假设的方式。
它既可以用于估计重构假设的期望风险，也可以在压缩规模受到统一控制时证明可学习性。

\subsection{从样本到假设的重构}

考虑递增阈值$h_t(x)=\mathbb I\{x\geq t\}$。四个输入0.1、0.3、0.7、0.9的标签
依次为0、0、1、1。只保留最小的正例$(0.7,1)$，并约定将保留输入作为阈值，
就能重构$h_{0.7}$。虽然其余三个观测已经被丢弃，这个假设仍然将四个输入全部分对。
这里重构的是一个与样本一致的假设，并不要求恢复生成数据的真实阈值。

这一办法适用于每份可由递增阈值正确分类的有限样本。只要样本中有正例，就保留最小的
正例；所有负例都位于它的左侧，其他正例都位于它的右侧或与它重合，因而重构阈值与
全部样本一致。若样本没有正例，则不保留任何观测，并约定由空样本重构处处预测0的假设。
这个假设可能不属于实数阈值空间，因此本节允许重构假设不属于原假设空间。

一般的压缩方案包含两个预先确定的步骤。压缩映射读取完整样本，选择少量带标签观测；
重构映射只能读取这些被保留的观测，并据此输出假设。有时还需要保留一个有限长度的
二进制串，以说明采用哪一种重构方式。这部分称为\term{旁信息}（side information），
它也必须计入压缩规模，不能借此保存任意多的原样本信息。

\begin{definition}[一致样本压缩方案]
\label{def:sample-compression}
设$\mathcal H\subseteq\{0,1\}^X$，$k,b$为非负整数。
一个$(k,b)$一致样本压缩方案由固定的确定性映射$\kappa$和$\rho$组成。
对任意有限样本$S=((x_i,y_i))_{i=1}^n$，若存在$h\in\mathcal H$使
$h(x_i)=y_i$对所有$i$成立，则$\kappa(S)=(S_I,u)$满足
\[
 I\subseteq\{1,\ldots,n\},\qquad |I|\leq k,\qquad u\in\{0,1\}^b,
\]
其中$S_I$按原样本中的次序排列。重构映射仅以$(S_I,u)$为输入，输出
$\rho(S_I,u):X\to\{0,1\}$，且满足
\[
 \rho(S_I,u)(x_i)=y_i,\qquad i=1,\ldots,n.
\]
\end{definition}

这个定义要求同一对映射对每份可实现样本都有效，且$k,b$不随样本量增长。
重构映射不能读取未保留的观测，也不能获得未计入旁信息的原始位置编号。
前面的阈值方案满足$k=1,b=0$；保留一个正例还是空样本，已经足以区分两种重构方式。

\subsection{期望风险与可学习性保证}

仅仅重构出一个与样本一致的假设，还不能直接说明其期望风险小。压缩规模的作用在于
限制算法能以多少种方式构造这个假设。对一份大小为$n$的样本，保留$j$个观测有
$\binom nj$种位置选择，旁信息有$2^b$种。因此，总共只需考虑
$2^b\sum_{j=0}^k\binom nj$种重构方式。

对于其中一种预先固定的选择，重构假设只依赖被保留的观测，其余观测仍然可以检验它的
预测是否正确。如果它的期望风险很大，却把这些独立观测全部分对，这种情形的概率就很小。
压缩算法虽然是在看过完整样本后才作出选择，但可以把每种选择对应的失败概率相加，
从而控制最终输出的期望风险。

\begin{theorem}[一致样本压缩的期望风险上界]
\label{thm:sample-compression-bound}
设$\mathcal H$具有一个$(k,b)$一致样本压缩方案，且压缩、重构及输出损失满足
第~\ref{sec:learning-problem}节的可测性与可分性约定。
若$S\sim\mathcal P^n$、$n>k$，且$\mathcal P$对$\mathcal H$可实现，
则对任意$\delta\in(0,1)$，以至少$1-\delta$的概率，重构假设
$h_S=\rho(\kappa(S))$的零一损失期望风险满足
\begin{equation}
 R_{\mathcal P}(h_S)
 \leq\frac{b\log2+\log\!\left(\sum_{j=0}^k\binom nj\right)+\log(1/\delta)}{n-k}.
 \label{eq:compression-generalization}
\end{equation}
\end{theorem}

\begin{proof}
固定位置集$I\subseteq\{1,\ldots,n\}$、$|I|=j\leq k$及旁信息$u$。
这里只固定一种可能的重构方式，不附加“压缩算法恰好选中了它”的条件。
给定$S_I$后，假设$h_{I,u}=\rho(S_I,u)$便已确定，其余$n-j$个观测仍然
相互独立且服从$\mathcal P$。若$R_{\mathcal P}(h_{I,u})>\varepsilon$，
它把这些观测全部分对的条件概率至多为
\[
 (1-\varepsilon)^{n-j}\leq e^{-(n-j)\varepsilon}
 \leq e^{-(n-k)\varepsilon}.
\]
再对$S_I$取平均，便知对这一固定的$(I,u)$，同时出现期望风险超过$\varepsilon$
和与全部样本一致的概率至多为$e^{-(n-k)\varepsilon}$。

压缩方案最终输出的假设必然对应某个$(I,u)$，并且与全部样本一致。
因此，若最终输出的期望风险超过$\varepsilon$，上述失败情形至少在一种选择中发生。
把所有选择对应的概率相加，得到
\[
 \mathbb P\{R_{\mathcal P}(h_S)>\varepsilon\}
 \leq 2^b\left(\sum_{j=0}^k\binom nj\right)e^{-(n-k)\varepsilon}.
\]
令右侧不超过$\delta$，对数变换后即得式~\eqref{eq:compression-generalization}。
若该式右侧大于或等于1，结论由零一损失期望风险不超过1直接成立。
\end{proof}

分子中的$b\log2$和组合数的对数，分别记录旁信息与保留位置的选择数；分母中的$n-k$则是至少还有多少个
未保留观测可以检验重构假设。这里的$\log$为自然对数，所以$b$比特对应$b\log2$。
对阈值方案，$k=1,b=0$，且$\sum_{j=0}^1\binom nj=n+1$，于是
\[
 R_{\mathcal P}(h_S)\leq\frac{\log(n+1)+\log(1/\delta)}{n-1}.
\]
例如$n=100$、$\delta=0.05$时，这个上界约为0.077：以至少95\%的概率，
重构阈值的期望错误率不超过7.7\%。这是通用压缩分析给出的充分保证，
不意味着阈值学习的最优样本量必须包含$\log n$因子。

\begin{theorem}[有界样本压缩推出PAC可学习]
\label{thm:compression-pac-learning}
在上述可测性条件下，若$\mathcal H$存在固定有限的$(k,b)$一致样本压缩方案，
则学习算法$A(S)=\rho(\kappa(S))$使$\mathcal H$可实现PAC可学习。
具体地，只要$n>k$且
\[
 b\log2+\log\!\left(\sum_{j=0}^k\binom nj\right)+\log(1/\delta)
 \leq(n-k)\varepsilon,
\]
就有$\mathbb P\{R_{\mathcal P}(A(S))\leq\varepsilon\}\geq1-\delta$，
且这一样本量条件不依赖具体的可实现分布$\mathcal P$。
\end{theorem}

\begin{proof}
当$k=0$时，组合数之和为1；当$k\geq1$固定时，
$\sum_{j=0}^k\binom nj\leq(k+1)n^k$。
因此，前一定理的分子至多按$\log n$的量级增长，分母则按$n$增长。
对任意$\varepsilon,\delta\in(0,1)$，都存在一个有限样本量阈值，使此后所有$n$
都满足所列条件。前一定理随即给出PAC要求的期望风险保证。
\end{proof}

这两个定理承担不同作用：前者估计某次学习输出的期望风险，后者说明同一个算法可以在
所有可实现分布上达到任意给定精度。它们直接依赖与样本一致的条件；存在标签噪声时，
重构假设未必能分对全部样本，不能直接套用这个证明。它们也只给出统计保证，
压缩与重构是否容易计算仍需单独分析。

\subsection{与VC维及Rademacher复杂度的关系}

VC维和样本压缩从两个方向描述假设空间。VC维考察一组输入能实现多少种不同预测，
样本压缩考察与样本一致的假设能够由多少个观测重构。在二分类设定中，有界的一致
压缩方案推出可实现PAC可学习，再由二分类学习基本定理推出VC维有限。
也可以直接通过计数看到这层限制。

\begin{theorem}[压缩规模对打散的限制]
\label{thm:compression-shattering-count}
若$\mathcal H$具有$(k,b)$一致样本压缩方案，并能打散$m$个不同输入，则
\[
 2^m\leq 2^b\sum_{j=0}^{\min\{k,m\}}\binom mj2^j.
\]
特别地，固定有限的$k,b$意味着$\operatorname{VC}(\mathcal H)<\infty$。
\end{theorem}

\begin{proof}
将这$m$个输入按固定次序排列。由于它们能被打散，全部$2^m$种0/1预测都对应一份
可实现的带标签样本。保留$j$个输入有$\binom mj$种选择，它们的标签有$2^j$种，
再加上$2^b$种旁信息，得到右侧的描述数。两个不同的完整预测不能被压缩成同一份
描述，否则固定的重构映射就必须对同一输入同时给出两个不同标签。
因此描述数至少为$2^m$。对固定$k,b$，右侧随$m$至多按多项式增长，
左侧按指数增长，不等式不可能对任意大的$m$成立。
\end{proof}

反过来，有限VC维也保证存在规模不依赖样本量的压缩方案。Moran与Yehudayoff证明，
VC维为$d$的二分类假设空间存在规模为$2^{O(d)}$的样本压缩方案
\scite{moran2016compression}。因此，在通常的可测性条件下，有限VC维、二分类PAC
可学习性和存在有界规模的压缩方案相互联系。不过，存在性并未说明能压缩到多小。
\term{样本压缩猜想}（sample compression conjecture）进一步要求与$d$成线性量级的
压缩规模；
Attias等人在2025年的研究中仍将这一一般性问题列为未解决问题
\scite{attias2025compression}。比较不同压缩规模时，需要同时核对保留观测数和旁信息量，
不能把大量描述信息藏在旁信息中。

Rademacher复杂度与样本压缩则采用了不同的证明对象。Rademacher复杂度控制整个
损失函数空间在样本上的随机波动，所得期望风险上界对空间中的所有假设同时成立；
压缩分析只要求控制由指定重构映射和有限描述产生的假设，不必证明整个假设空间都满足
同样精度的一致收敛。前者衡量函数值适应随机符号的能力，后者计数可能的样本描述，
所以压缩规模不能直接作为Rademacher复杂度代入公式。

两种方法都需要防止算法利用样本中的偶然性：Rademacher分析对整个函数空间控制这种
选择的影响，压缩分析则限制重构假设所依赖的信息。上一节的有界损失上界允许非零经验风险；
本节的一致压缩上界利用了经验风险为零，通常给出固定$k,b$下约为$\log n/n$的衰减量级。
两者的条件不同，不能仅凭这个量级比较优劣。即使某次样本碰巧容易压缩，也不自动说明
整个假设空间可学习；PAC结论需要一个对所有可实现样本有效、规模统一受控的方案。

\section{集成学习}
\label{sec:boosting}

前面讨论了如何评估一个假设的期望风险。本节进一步考虑，能否把多个假设组合起来，
得到比单个假设更可靠的预测。\term{集成学习}（ensemble learning）将多个基假设的输出
通过投票或平均形成最终预测。组合既可能降低单次训练的随机波动，也可能表达单个
基假设无法表达的分类边界；但增加假设数本身并不保证期望风险下降。

本节围绕两个问题展开：投票在什么条件下能够改善预测，以及如何主动构造满足要求的
基假设。前者需要分析错误之间的关系，后者由提升方法给出一种具体实现。
最终还要控制投票假设的期望风险，才能把样本上的改善转化为可学习性结论。

\subsection{投票、平均与错误相关性}

本节将标签写成$\{-1,1\}$，并约定得分为0时预测1。给定基假设
$h_1,\ldots,h_T:X\to\{-1,1\}$，以及非负权重$a_t$、$\sum_ta_t=1$，定义
\[
 f(x)=\sum_{t=1}^T a_t h_t(x),\qquad H(x)=\operatorname{sign}(f(x)).
\]
$f(x)$是加权投票得分，$H(x)$是最终分类假设。取$a_t=1/T$便得到多数投票。
若允许从基假设空间$\mathcal B$选取任意有限个假设，则这些得分函数构成
$\mathcal B$的\term{凸包}（convex hull），记为$\operatorname{conv}(\mathcal B)$。
这里凸包中的元素是实值得分函数，不能与阈值化之后的二分类假设混为一谈。

三个假设在某个输入上的预测为$(1,1,-1)$时，多数投票输出1。
若真实标签是1，这次投票纠正了第三个假设的错误；若三个假设都预测为$-1$，
则投票无法纠正它们共同的错误。因此，需要关注基假设是否在同一批输入上犯错。
下面的独立性条件给出一个理想化的充分保证。

\begin{theorem}[独立错误下的多数投票]
\label{thm:independent-majority-vote}
设$T$为奇数。对$\mathcal P$几乎处处的$(x,y)$，假设由训练随机性产生的错误指标
$E_t=\mathbb I\{h_t(x)\neq y\}$在给定$(x,y)$后相互独立，且
$\mathbb E[E_t\mid x,y]\leq1/2-\gamma$，其中$\gamma>0$。
则等权多数投票满足
\[
 \mathbb E_{\mathrm{train}}R_{\mathcal P}(H)\leq e^{-2T\gamma^2}.
\]
\end{theorem}

\begin{proof}
给定$(x,y)$，投票出错意味着$T^{-1}\sum_tE_t\geq1/2$。
这一平均值的条件期望至多为$1/2-\gamma$，由独立有界变量的Hoeffding不等式，
出错概率至多为$e^{-2T\gamma^2}$。再对$(x,y)\sim\mathcal P$取期望即得结论。
\end{proof}

这里要求每个输入处的条件错误率受控，而且错误条件独立；仅仅知道每个基假设的
总体期望风险小于$1/2$并不足够。反复训练同一个算法时，难分类的输入可能让所有
假设一起出错。这个定理解释了相互独立的错误为何有利于投票，但不能直接用作
一般集成算法的期望风险保证。

\term{装袋法}（bootstrap aggregating, bagging）通过对原样本有放回地重采样，
训练多个基假设，再平均得分或投票。它主要尝试减小训练样本扰动带来的预测变化。
例如，在固定输入$x$处，若实值得分$F_1(x),\ldots,F_T(x)$各有方差$v$，
两两协方差均为$c$，则展开平均值的方差可得
\[
 \operatorname{Var}\!\left(\frac1T\sum_tF_t(x)\right)
 =\frac{v}{T}+\frac{T-1}{T}c.
\]
当$c=0$时，平均得分的方差为$v/T$；当$c=v$时，平均并不降低方差。
这里的随机性来自训练过程和数据，重采样共享同一份原样本，不能据此假定这些得分
无条件独立。方差降低也不自动等于零一损失期望风险降低，因为分类还取决于得分
是否越过0。随机森林在重采样之外随机选择分裂特征，也是在构造有差异的基假设；
其具体算法留待模型章节讨论\scite{shalevshwartz2014}。

\subsection{弱学习与提升}

第~\ref{chap:statistical-learning-framework}章把弱学习描述为获得一个固定优势：在每个
可实现分布上，统一算法都能以高概率输出期望风险不超过$1/2-\gamma$的分类器，其中
$\gamma\in(0,1/2)$与分布无关。一次这样的学习只保证优于随机猜测，尚未达到任意小的
误差容限。\term{提升}（boosting）反复调用弱学习器，并改变每一轮关注的样本，
使多个弱分类器的错误相互补偿。

弱学习要求同一个学习器对每个允许的可实现分布都获得优势，而不只是对原始分布有效。
提升算法会改变各个样本的权重，这相当于不断提出新的分类分布；若弱学习器只能在
最初的分布上优于随机猜测，就不足以支持后续保证。这里的弱优势可以很小，
但必须在这些分布上统一成立。

\subsection{经验风险}

AdaBoost是提升方法的一个具体构造。它给训练样本赋予非负权重；某轮分错的样本在下一轮
获得更大权重。每轮学到的假设还会参与最终加权投票。经验风险的下降可以直接由权重更新
证明，但要得到PAC保证，还需单独控制最终投票的期望风险。

沿用本节的$-1,1$标签和基假设输出约定，
$y_i h_t(x_i)=1$表示分对，$y_i h_t(x_i)=-1$表示分错。令初始样本权重
$D^{(1)}(i)=1/n$，其中$D^{(t)}(i)$是第$t$轮分配给样本$i$的权重，且$\sum_iD^{(t)}(i)=1$。
弱学习器返回假设$h_t$后，算法计算其加权错误率
$e^{(t)}=\sum_iD^{(t)}(i)\mathbb I\{h_t(x_i)\neq y_i\}$。
在$0<e^{(t)}<1/2$时，令
\begin{equation}
 \alpha^{(t)}=\frac12\log\frac{1-e^{(t)}}{e^{(t)}},\qquad
 D^{(t+1)}(i)=\frac{D^{(t)}(i)e^{-\alpha^{(t)} y_i h_t(x_i)}}{Z^{(t)}},
 \label{eq:adaboost-update}
\end{equation}
其中$Z^{(t)}=\sum_iD^{(t)}(i)e^{-\alpha^{(t)} y_i h_t(x_i)}$使新权重和为1。
若$e^{(t)}=0$，当前假设已与全部样本一致，可直接输出该假设并停止；若$e^{(t)}\geq1/2$，本轮不满足所需优势。
在预先固定平局规则后，最终输出$H_T(x)=\operatorname{sign}(\sum_t\alpha^{(t)} h_t(x))$。

重加权改变的是每个样本在下一轮错误率中所占的权重。若一个
样本被当前弱分类器分错，指数因子为$e^{\alpha^{(t)}}$，其相对权重上升；分对时因子为
$e^{-\alpha^{(t)}}$，相对权重下降。下一轮弱学习器若仍想取得小于$1/2$的加权错误，就必须
更多照顾此前失败的样本。$e^{(t)}$越小，$\alpha^{(t)}$越大，最终投票中该轮假设的权重也越大。

一次数值更新可以说明公式的作用。设四个样本的初始权重都是$1/4$，本轮只分错第一个，
则$e^{(t)}=1/4$、$\alpha^{(t)}=\tfrac12\log3$。归一化后，第一个样本的权重变为$1/2$，
另外三个各为$1/6$。下一轮若继续把第一个样本分错，仅这一次错误就占权重的一半，
弱学习器必须改变判断才能保持严格优于$1/2$的加权错误率。

\begin{theorem}[AdaBoost经验风险界]
\label{thm:adaboost-training-error}
若每轮的加权错误率满足$0<e^{(t)}<1/2$，并记$\gamma^{(t)}=1/2-e^{(t)}$，则
\begin{equation}
 \hat R_S(H_T)\leq\prod_{t=1}^T Z^{(t)}
 =\prod_{t=1}^T2\sqrt{e^{(t)}(1-e^{(t)})}
 \leq\exp\left(-2\sum_{t=1}^T(\gamma^{(t)})^2\right).
 \label{eq:adaboost-training-bound}
\end{equation}
\end{theorem}

\begin{proof}
递推式展开后有
$D^{(T+1)}(i)=n^{-1}\exp(-y_i\sum_t\alpha^{(t)} h_t(x_i))/\prod_tZ^{(t)}$。
若$H_T$在样本$i$上出错，则$y_i\sum_t\alpha^{(t)} h_t(x_i)\leq0$，指数项至少为$1$。
将上式对所有训练样本求和，并利用$\sum_iD^{(T+1)}(i)=1$，得到误分类比例不超过
$\prod_tZ^{(t)}$。
按正确与错误两组求和可得$Z^{(t)}=2\sqrt{e^{(t)}(1-e^{(t)})}$，最后用
$\sqrt{1-4(\gamma^{(t)})^2}\leq e^{-2(\gamma^{(t)})^2}$。
\end{proof}

\subsection{可学习性}

若每轮都有$e^{(t)}\leq1/2-\gamma$，取
$T\geq\log(2/\varepsilon)/(2\gamma^2)$，定理~\ref{thm:adaboost-training-error}
便使经验风险不超过$\varepsilon/2$。要进一步保证期望风险小，还要控制算法根据样本
选择基假设和投票权重所带来的偏差。这里并不需要基假设空间有限；
关键在于它在有限输入上能够给出多少种不同预测。

例如，实数阈值有无限多个，但在$m$个有序输入上只能给出$m+1$种不同预测。
选择$T$个阈值时，即使每个阈值的数值都有无限种可能，在这批输入上产生的
$T$列预测结果也至多有$(m+1)^T$种。一般的无限基假设空间可以用生长函数作同样的计数。

设$\mathcal B\subseteq\{-1,1\}^X$是抽样前固定的基假设空间，$\mathcal V_T$
表示由$T$个基假设和任意实数权重形成的投票分类假设空间，平局规则固定。
将$-1,1$重新编码为0,1不改变预测数目，因此仍沿用前文的生长函数及VC维记号。
以下讨论沿用第~\ref{sec:learning-problem}节的可测性与可分性约定。

\begin{theorem}[投票假设空间的生长函数]
\label{thm:vote-growth-infinite}
对任意$T\geq1$及$m\geq T+1$，有
\begin{equation}
 \Pi_{\mathcal V_T}(m)
 \leq\Pi_{\mathcal B}(m)^T
       \left(\frac{em}{T+1}\right)^{T+1}.
 \label{eq:boosting-vote-growth}
\end{equation}
若$d=\operatorname{VC}(\mathcal B)$满足$1\leq d<\infty$，则在
$m\geq\max\{d,T+1\}$时，
\begin{equation}
 \log\Pi_{\mathcal V_T}(m)
 \leq dT\log\frac{em}{d}+(T+1)\log\frac{em}{T+1}.
 \label{eq:infinite-boosting-growth}
\end{equation}
特别地，对每个固定有限的$T$，$\mathcal V_T$的VC维有限。
\end{theorem}

\begin{proof}
固定$m$个输入。每个基假设在这些输入上有至多$\Pi_{\mathcal B}(m)$种不同预测，
所以$T$个基假设形成的预测矩阵至多有$\Pi_{\mathcal B}(m)^T$种。
固定其中一个矩阵后，每个输入对应一个$T$维向量
$(h_1(x),\ldots,h_T(x))$；改变投票权重，就是在这些向量上改变一个线性分类器。
这个分类器属于$T$维仿射半空间假设空间，其VC维为$T+1$。
Sauer--Shelah引理因此将每个固定矩阵上的投票预测数目限制为
$(em/(T+1))^{T+1}$。对所有矩阵计数，得到第一式。

当基假设空间的VC维为$d$时，再代入
$\Pi_{\mathcal B}(m)\leq(em/d)^d$，便得到第二式。
固定$d,T$后，这个上界对$\Pi_{\mathcal V_T}(m)$只允许多项式增长，
因此不可能对任意大的$m$都达到打散所需的$2^m$，故投票空间的VC维有限。
若$d=0$，非空$\mathcal B$在每个输入上的预测唯一，第一式同样给出有限VC维。
\end{proof}

这个证明将无限空间的计数转化为有限输入上的预测计数。
若$\mathcal B$本来就有限，还可使用$\Pi_{\mathcal B}(m)\leq|\mathcal B|$，
得到熟悉的$T\log|\mathcal B|$项。因此，有限基假设空间只是上述分析的一种特殊情形。
对无限但VC维有限的基假设空间，每个固定轮数的投票空间同样满足一致收敛。

\begin{theorem}[有限VC维基假设空间下的弱学习提升]
\label{thm:finite-class-boosting}
\label{thm:finite-vc-boosting}
设$\mathcal H$为目标假设空间，$\mathcal B$为固定的非空基假设空间，
$\operatorname{VC}(\mathcal B)<\infty$，允许$|\mathcal B|=\infty$。
假设存在弱学习器，对每个$\mathcal H$可实现分布和每个$\eta\in(0,1)$，
使用统一有限数量的独立样本后，以至少$1-\eta$的概率输出$h\in\mathcal B$，
其零一损失期望风险不超过$1/2-\gamma$，其中$\gamma>0$与分布无关。
则允许从原样本的加权经验分布重新抽样的AdaBoost构造，能够非恰当地可实现PAC学习
$\mathcal H$。
\end{theorem}

\begin{proof}
给定$\varepsilon,\delta$，在抽取原样本前固定
$T=\lceil\log(2/\varepsilon)/(2\gamma^2)\rceil$。
由定理~\ref{thm:vote-growth-infinite}，$\mathcal V_T$的VC维有限，故可以选择
与具体分布无关的原始样本量$n$，使其以至少$1-\delta/2$的概率满足
\[
 \sup_{H\in\mathcal V_T}|R_{\mathcal P}(H)-\hat R_S(H)|\leq\varepsilon/2.
\]

原始分布对$\mathcal H$可实现，所以原样本几乎必然与某个目标假设一致。
每轮改变观测权重并不改变标签，因此加权经验分布也对$\mathcal H$可实现。
条件于原样本和此前各轮结果，从当前加权经验分布独立重采样，再以失败概率
$\delta/(2T)$调用弱学习器，即能以相应概率获得加权错误率不超过$1/2-\gamma$的
基假设。把各轮失败概率相加，所有轮次都获得优势的概率至少为$1-\delta/2$。

在两类成功事件同时发生时，AdaBoost的经验风险不超过$\varepsilon/2$，
而其输出属于$\mathcal V_T$，所以期望风险不超过$\varepsilon$。
若提前找到经验风险为零的基假设，则直接输出它；其余投票权重取零即可将其视为
$\mathcal V_T$中的元素。总失败概率至多为$\delta$，其中包含原样本抽样和各轮学习的
随机性，故得到PAC保证。
\end{proof}

因此，无限基假设空间并不妨碍提升；有限VC维和对所有可实现分布统一成立的弱优势，
足以支撑这里的构造。但有限VC维本身并不保证基假设能够弱学习任意给定的目标空间
$\mathcal H$，弱学习条件仍须另外满足。这一定理也没有假定各轮错误相互独立
\scite{shalevshwartz2014}。

上述证明对给定精度先固定$T$，再选择足够的样本量。
若基假设空间的VC维无限，或者在固定样本量下任意增加轮数，便不能直接依靠这个
有限VC维结论获得统一保证。即使每轮经验风险继续改善，也还需要其他复杂度条件来
控制期望风险。存在标签噪声时，每个重加权分布上的弱优势也可能不再成立，
因此这里的可实现证明不能直接扩展为不可知提升保证。

\subsection{估计期望风险}

按投票轮数计数可以证明可学习性，但可能高估一个实际集成模型的复杂度。
另一种分析利用上一节的Rademacher复杂度，研究归一化投票得分
$f=\sum_ta_th_t$。在样本$(x,y)$上，$yf(x)$表示正确一侧的总权重减去错误一侧的
总权重。例如，正确一侧占0.9时，间隔为0.8；只占0.51时，间隔为0.02。
两者都分类正确，但后一种投票更接近改变类别的边界。

\begin{theorem}[凸组合的Rademacher复杂度]
\label{thm:learning-ensemble-convex-hull}
对任意固定输入样本$S_X$和非空有界实值函数空间$\mathcal B$，有
\[
 \widehat{\mathfrak R}_{S_X}(\operatorname{conv}(\mathcal B))
 =\widehat{\mathfrak R}_{S_X}(\mathcal B).
\]
\end{theorem}

\begin{proof}
固定随机符号后，加权和是函数值的线性函数。对任意凸组合，
\[
 \frac1n\sum_i\sigma_i\sum_ta_th_t(x_i)
 =\sum_ta_t\left(\frac1n\sum_i\sigma_ih_t(x_i)\right)
 \leq\sup_{h\in\mathcal B}\frac1n\sum_i\sigma_ih(x_i).
\]
而每个基假设本身也属于凸包，两边的上确界相等。对随机符号取期望即得结论。
\end{proof}

取上一节定义的间隔损失$\ell_\rho$，其中$\rho>0$在抽样前固定。
它不小于零一损失，而且关于得分的Lipschitz常数为$1/\rho$。
对损失空间应用收缩引理及经验Rademacher期望风险上界，便得到：以至少$1-\delta$
的概率，所有$f\in\operatorname{conv}(\mathcal B)$同时满足
\begin{equation}
 R_{\mathcal P}^{01}(\operatorname{sign}f)
 \leq\frac1n\sum_i\ell_\rho(f(x_i),y_i)
 +\frac{2}{\rho}\widehat{\mathfrak R}_{S_X}(\mathcal B)
 +3\sqrt{\frac{\log(2/\delta)}{2n}}.
 \label{eq:ensemble-margin-bound}
\end{equation}
对固定标签乘上$y_i$不会改变随机符号的分布，所以这里使用输入样本上的基假设复杂度。
这一推导也说明，凸包复杂度的等式控制的是实值得分，并非直接控制阈值化后的零一损失。

这个上界同样适用于无限基假设空间，因为凸包复杂度的证明并未使用基假设总数。
例如，若$\mathcal B\subseteq\{-1,1\}^X$的VC维为$d$，则在$1\leq d\leq n$时，
对固定样本上的有限预测向量应用有限类Rademacher上界，再用Sauer--Shelah引理，可得
\[
 \widehat{\mathfrak R}_{S_X}(\mathcal B)
 \leq\sqrt{\frac{2\log\Pi_{\mathcal B}(n)}{n}}
 \leq\sqrt{\frac{2d\log(en/d)}{n}}.
\]
这里每个预测向量的欧氏范数为$\sqrt n$，所以即使基假设有无限多个，
复杂度仍由有限样本上的预测数目控制。代入式~\eqref{eq:ensemble-margin-bound}，
复杂度项至多为$\frac{2}{\rho}\sqrt{2d\log(en/d)/n}$。
若VC维无限，这一替换不能使用，但原始的经验Rademacher上界仍然有效；
此时需要实际得到较小的复杂度，才能给出有用的期望风险保证。

式~\eqref{eq:ensemble-margin-bound}没有显式出现基假设的个数$T$。
因此，在基假设空间固定时，增加轮次若能改善样本上的归一化间隔，可能使这个期望风险
上界减小。它并不保证增加轮次一定改善间隔，也不保证所有集成模型都优于单个模型。
集成学习的统计保证需要同时检查组合方式、基假设的复杂度，以及样本上的错误或间隔；
它的计算代价还取决于训练轮数和每个基假设的训练、预测时间。

\section{可计算PAC学习}
\label{sec:computable-binary-learning}

前面的统计结论说明，在样本足够多时，某种学习规则能够输出期望风险较小的假设。
但是，$\argmin_{h\in\mathcal H}\hat R_S(h)$只是指定了要找什么，并没有说明如何找到它。
同样，VC维有限并不自动提供寻找假设的程序，压缩方案的存在也不保证压缩能够快速完成。
要把统计保证变成可运行的学习算法，需要分别考察程序是否能够完成计算，以及完成计算
需要多少资源。

\subsection{统计保证与程序实现}

讨论计算之前，必须说明数据与假设怎样交给计算机。例如，布尔向量可用有限比特串
表示，自然数可用二进制表示；实数输入则需要另行约定有理数编码、精度或访问模型。
输入维数相同而数值编码长度不同，计算代价也可能不同，因此不能只用样本个数衡量运行时间。
本节关于可计算性的结论使用$X=\mathbb N$，假设以能对每个输入停机并输出0或1的
程序表示。关于高效学习的例子则使用有限维布尔输入。

\begin{definition}[可计算PAC学习]
\label{def:binary-computable-pac}
设$\mathcal H\subseteq\{0,1\}^{\mathbb N}$。若存在一个对每份有限带标签样本$S$
都停机的程序$A$，输出一个总可计算二分类假设的程序编码，并且对任意
$\varepsilon,\delta\in(0,1)$存在有限的$m(\varepsilon,\delta)$，使所有
$n\geq m(\varepsilon,\delta)$及所有分布$\mathcal P$均满足
\[
 \mathbb P_{S\sim\mathcal P^n}\left\{
 R_{\mathcal P}(A(S))\leq\inf_{h\in\mathcal H}R_{\mathcal P}(h)+\varepsilon
 \right\}\geq1-\delta,
\]
则称$\mathcal H$是不可知\term{可计算PAC可学习的}（computably PAC learnable, CPAC）。
将分布限制为$\mathcal H$可实现分布，并将右侧风险要求改为
$R_{\mathcal P}(A(S))\leq\varepsilon$，得到可实现版本。
\end{definition}

“总可计算”意味着预测程序对每个自然数输入都停机。若还要求$A(S)\in\mathcal H$，
则是恰当学习；允许输出其他可计算假设，则是非恰当学习。
定义要求存在足够大的样本量，但没有要求程序能根据$\varepsilon,\delta$算出一个
足够的样本量。当样本量上界本身也要求可计算时，文献称之为strong CPAC。
这一术语中的strong指样本量上界的有效性，与集成学习中将弱优势提升到任意精度的
“强学习”含义不同。

对于有限且给定可执行列表的假设空间，可以逐个计算经验风险，再输出经验风险最小者。
问题出现在无限假设空间：即使能够依次列出所有假设，也未必知道何时已经检查完所有
可能更好的假设。特别是在有噪声时，当前最小经验风险仍可能被后续假设改进。
因此，“可以枚举假设”和“可以计算ERM”是不同条件。

\begin{theorem}[可计算ERM给出的学习保证]
\label{thm:computable-erm-pac}
设$\mathcal H\subseteq\{0,1\}^{\mathbb N}$的VC维$d$有限。
若存在对所有有限带标签样本都停机的程序，输出$\mathcal H$中一个可有效求值的
经验风险最小化假设，则$\mathcal H$恰当地不可知CPAC可学习，且具有可计算的样本量上界。
\end{theorem}

\begin{proof}
令$A$为给定ERM程序。有限VC维给出分布无关的一致收敛样本量上界；
将精度设为$\varepsilon/2$，在一致收敛事件上，对任意$h\in\mathcal H$有
\[
 R_{\mathcal P}(A(S))
 \leq\hat R_S(A(S))+\varepsilon/2
 \leq\hat R_S(h)+\varepsilon/2
 \leq R_{\mathcal P}(h)+\varepsilon.
\]
对$h$取下确界得到所需保证。VC一致收敛上界是$d,\varepsilon,\delta$的显式函数，
所以对固定假设空间，将一个有限的VC维上界作为常数写入程序，就能计算足够的样本量。
训练程序与输出假设的可计算性则由题设保证。
\end{proof}

证明中使用已知的有限维度上界，并不意味着能从任意假设空间的程序描述中自动算出VC维。
这个区别体现了统计存在性与统一算法构造之间的差距。

\subsection{有效VC维}

普通VC维为$d$意味着：任取$d+1$个不同输入，总有一种0/1预测不能由假设空间实现。
这个结论只断言这样的预测存在。对计算机而言，还需要一套步骤，能够读入这些输入，
在有限时间内找出一组不能实现的预测。\term{有效VC维}（effective VC dimension）
把这一寻找过程加入维度条件。

\begin{definition}[有效VC维]
\label{def:effective-vc-dimension}
对$\mathcal H\subseteq\{0,1\}^{\mathbb N}$，若存在总可计算映射
$W$，使任意$d+1$个互异输入$x_1,\ldots,x_{d+1}$都满足
\[
 W(x_1,\ldots,x_{d+1})\in\{0,1\}^{d+1}
 \setminus\{(h(x_1),\ldots,h(x_{d+1})):h\in\mathcal H\},
\]
则称$d$为有效VC维的一个上界。所有这样的非负整数$d$中的最小值记为
$\operatorname{eVC}(\mathcal H)$；若不存在，则记为$\infty$。
\end{definition}

映射$W$只需给出一种不可能的预测，无须列出所有可能预测。
对自然数上的递增阈值，给定两个不同输入，比较大小后令较小者对应1、较大者对应0，
就得到无法实现的预测。这个程序总能停机，因此有效VC维至多为1。
单个输入又可分别预测为0或1，所以普通VC维和有效VC维都为1。

\begin{theorem}[普通VC维与有效VC维的比较]
\label{thm:effective-vc-comparison}
对非空$\mathcal H\subseteq\{0,1\}^{\mathbb N}$，
$\operatorname{VC}(\mathcal H)\leq\operatorname{eVC}(\mathcal H)$。
若普通VC维有限，且$\mathcal H$具有定理~\ref{thm:computable-erm-pac}所述的可计算ERM，
则两者相等。
\end{theorem}

\begin{proof}
若有效VC维至多为$d$，见证程序对任意$d+1$个输入都会给出一个不能实现的预测，
所以这些输入不可能被打散，普通VC维至多为$d$。

再设普通VC维为$d<\infty$。给定$d+1$个不同输入，枚举全部$2^{d+1}$种0/1赋值，
对每一份相应带标签样本运行ERM并计算其经验风险。每次计算都停机，且至少有一种赋值
不能被任何$h\in\mathcal H$完全实现。该赋值的最小经验风险大于零，因此ERM输出
也具有正经验风险。返回首次发现的这种赋值，就构造了有效见证程序。
故有效VC维至多为$d$，结合前一方向得证。
\end{proof}

有效VC维还给出非恰当可计算学习的完整刻画。在上述自然数输入、以总可计算函数为
输出的模型中，Delle Rose等证明：非恰当不可知CPAC可学习，当且仅当有效VC维有限
\scite{dellerose2023}。这一结果需要构造适当的可计算学习器，不能仅由普通VC维的
一致收敛证明推出。这里给出刻画以说明有效维度的作用；其完整计算论证明见原文。

允许非恰当输出是这条刻画的条件。若额外要求输出始终属于原假设空间，就不能直接
沿用结论；样本量上界是否可计算也可能产生区别。
有关研究构造了普通VC维有限、却不存在非恰当CPAC学习器的假设空间，说明
即使统计上的样本数量不是障碍，程序实现仍可能成为障碍
\scite{dellerose2023}。这类不可计算性也不同于下一小节的运行时间问题。

\subsection{高效学习与优化误差}

一个程序最终停机，仍可能需要无法承受的时间。高效PAC学习在可计算性的基础上，
要求样本量、训练时间和预测时间受多项式控制。讨论一族问题时，要明确多项式的变量：
通常包括输入维数或问题规模$p$、输入和输出的编码长度、$1/\varepsilon$以及
$\log(1/\delta)$；具体表示模型还可能引入目标假设的描述长度。
把一次实数运算当成一步，与逐次计算比特操作的数量，是不同的计算模型，不能混用。

统计分析也不一定要求精确求解ERM。允许算法输出经验风险只比最优值大$\alpha$的假设，
可以把求解精度与样本估计精度分别控制。

\begin{theorem}[近似ERM的期望风险保证]
\label{thm:approximate-erm-risk}
设$A(S)\in\mathcal H$，且
$\hat R_S(A(S))\leq\inf_{h\in\mathcal H}\hat R_S(h)+\alpha$。
若样本满足$\sup_{h\in\mathcal H}|R_{\mathcal P}(h)-\hat R_S(h)|\leq\eta$，则
\[
 R_{\mathcal P}(A(S))
 \leq\inf_{h\in\mathcal H}R_{\mathcal P}(h)+2\eta+\alpha.
\]
\end{theorem}

\begin{proof}
对任意$h\in\mathcal H$，依次使用一致收敛、近似最优性及一致收敛，得到
\[
 R_{\mathcal P}(A(S))
 \leq\hat R_S(A(S))+\eta
 \leq\hat R_S(h)+\alpha+\eta
 \leq R_{\mathcal P}(h)+\alpha+2\eta.
\]
对$h$取下确界即得结论。
\end{proof}

这里$2\eta$来自经验风险与期望风险的偏差，$\alpha$来自优化没有精确求解。
例如取$\eta=\varepsilon/4$、$\alpha=\varepsilon/2$，就能将总的超额期望风险
控制在$\varepsilon$以内。若获得这种近似解的时间和所需样本量都是所约定参数的
多项式，就得到高效不可知PAC学习的一条充分途径。
相反，某个精确ERM问题难求解，尚不足以证明所有近似或非恰当学习器都不高效。

\subsection{从枚举到有效算法}

布尔合取展示了统计复杂度与计算方法的区别。假设每个输入含$p$项是非特征，例如邮件是否含附件、是否出现某个词，记为$\bm x\in\{0,1\}^p$。一个条件可以要求$x_j=1$，也可以要求$x_j=0$；它们分别用$x_j$及其否定$1-x_j$表示，称为文字。合取假设要求所选条件全部满足，也就是所有选中文字均为1时预测1，否则预测0。
每个坐标可选择正文字、负文字或不使用，所以不含矛盾文字的合取至多有$3^p$个。
再加入恒为0的假设，得到有限假设空间$\mathcal C_p$。
有限假设空间的可实现样本量界因而只按$p$线性增长，但逐个枚举这些假设需要检查
指数数量的候选。枚举很慢只是这种求解方法的性质。

在可实现设定下，可以直接构造与样本一致的合取：初始保留全部$2p$个文字，
每遇到一个标签为1的样本，就删除在它上面取值为0的文字。处理完所有正例后，
输出剩余文字的合取；没有剩余文字时输出恒为1。若没有正例，则输出恒为0。

\begin{theorem}[布尔合取的高效可实现学习]
\label{thm:efficient-conjunction-learning}
上述算法在$\mathcal C_p$可实现的样本上输出与样本一致的假设，训练时间为$O(np)$，
预测时间为$O(p)$。对于$\varepsilon,\delta\in(0,1)$，只要
\[
 n\geq\frac{\log(3^p+1)+\log(1/\delta)}{\varepsilon},
\]
就以至少$1-\delta$的概率输出零一损失期望风险不超过$\varepsilon$的假设。
\end{theorem}

\begin{proof}
若样本没有正例，恒为0与样本一致。否则，目标合取中的每个文字都在所有正例上为1，
所以不会被删除。算法保留的每个文字也都在全部正例上为1，因此输出合取分对全部正例。
对于任意负例，目标合取中至少有一个文字在该输入上为0；这个文字仍然被保留，
故输出也预测0。算法因此与全部样本一致。每个样本最多检查$2p$个文字，
训练时间为$O(np)$，输出最多含$p$个非矛盾文字，预测时间为$O(p)$。
将$|\mathcal C_p|\leq3^p+1$代入有限假设空间的可实现PAC样本量界即得风险保证。
\end{proof}

这个例子同时给出了样本量、实际算法和运行时间，而不是只写出一个经验风险最小化目标。
它也依赖可实现条件：有噪声时，错误标成正例的观测可能删除目标合取需要的文字，
上述一致性证明就失效。改变学习设定后，需要重新分析优化与统计保证。

\section{本章小结}
\label{sec:binary-learning-summary}

要判断有限样本能否排除期望风险较高的候选，不能只分析一个事先固定的假设，还必须
控制算法看过样本后可能选择的整个假设空间。对有限假设空间，把所有假设的失败概率
相加，可以得到可实现与不可知设定下的充分样本量；假设空间无限时，计数对象改为它在
有限输入上给出的不同预测。
Sauer--Shelah引理用VC维控制这个预测数目，再经对称化建立经验风险与期望风险的联系。
在共享可测性约定及本章列明的条件下，有限VC维、一致收敛、可实现PAC可学习与不可知PAC可学习等价。

可学习性刻画回答是否存在统一的样本量保证，具体期望风险上界还可以利用更多信息。
Rademacher复杂度与高斯复杂度分别用随机正负号和高斯变量衡量函数空间拟合随机波动的能力，
样本压缩则利用重构输出所需保留的信息量。集成学习需要同时控制基假设的表达能力、组合方式与训练过程；
无限基假设空间并不妨碍学习，但经验风险随轮数下降也不能单独保证期望风险下降。
这些分析给出了不同的充分条件，不能只凭表达式较短就认定哪一种总是更紧。

统计维度有限、程序能够实现、运行时间为多项式，是三个不同层次。有效VC维与可计算
学习器联系起来，高效学习还需核算训练、预测和输出表示的代价。二分类中的这些区别在
多分类中仍然存在，但单一维度同时刻画一致收敛与可学习性的等价关系将不再普遍成立。

\paragraph{延伸阅读}
若希望进一步推导依赖样本的期望风险上界，可阅读Bartlett与Mendelson关于Rademacher
复杂度和高斯复杂度的论文。它不仅给出概率保证，还讨论组合函数空间的复杂度怎样由
基函数空间控制，适合接着本章的线性预测与集成学习分析阅读\scite{bartlettmendelson2002}。

对样本压缩感兴趣，可以继续阅读Moran与Yehudayoff关于VC假设空间压缩方案的构造。
关注点应放在“有限VC维为何能保证存在压缩表示”，并与本章“已有压缩表示怎样推出
期望风险上界”的证明对照。压缩表示的存在、大小以及能否高效求出仍需分别判断\scite{moran2016compression}。

可计算性方向可继续阅读Delle Rose等的有效VC维刻画，比较普通PAC、CPAC与附加可计算
样本量要求的版本。这有助于理解为什么统计学习基本定理不能直接充当可执行算法\scite{dellerose2023}。
